% !TeX root = ../main.tex

\xchapter{面向视频摘要的多智能体大语言模型推理框架}{Multi-agent large language model inference framework for video summarization}

上一章提出的基于层级化语义建模的免训练视频摘要方法，借助帧、场景和视频三级语义表示与大语言模型的零样本推理能力，在免训练条件下取得了具有竞争力的摘要性能。然而，该方法在重要性评估阶段仍依赖单一大语言模型对各场景逐一统一打分，存在三方面结构性瓶颈：单次推理难以充分兼顾叙事推进、视觉显著性、情绪表达与信息增益等多维评价目标；固定评估流水线无法针对不同视频动态调整评价策略；场景间独立评分缺乏对前序上下文的显式建模，在长视频中易引发局部评分不一致与冗余保留。针对上述局限，本章提出面向视频摘要的多智能体大语言模型推理框架，将单一模型的统一评分拆解为 Planner Agent 全局规划、多个 Expert Agent 多维协同评估与上下文记忆增强三位一体的结构化推理流程，实现从单源评分到多角色协作决策的范式跃迁。本章首先分析上述问题的成因与建模思路，继而阐述多智能体协同视频摘要方法的整体框架与各模块设计，随后介绍数据集、评价指标与实验设置，并通过对比实验与消融实验验证所提方法的有效性，最后对本章工作加以总结。

\xsection{问题提出}{Problem Statement}

第三章以摆脱监督训练与人工标注依赖为目标，提出了基于层级化语义建模的免训练视频摘要方法，通过构建帧、场景和视频三级语义表示并借助大语言模型进行场景重要性推理，在无需任务特定训练的条件下取得了具有竞争力的性能。然而，从推理架构层面审视该方法，可以发现其仍存在若干结构性局限。这些局限并非源于标注数据的匮乏或视觉特征的不足，而是根植于单一大语言模型与固定评估流水线相结合的范式在面对开放域视频时固有的表达瓶颈。

具体而言，上述瓶颈可从评价维度、评估策略与时序建模三个层面加以剖析。在评价维度层面，第三章方法将全部评价意图压缩到单次大语言模型调用中，由同一提示同时承载叙事推进、视觉显著性、情绪表达与信息密度等多维判断需求。然而，大语言模型在单次生成中能够有效关注的语义维度有限，当多个相互独立乃至彼此竞争的评价目标被同时施加时，模型输出不可避免地沦为各维度间的隐式折中，难以对任一维度给出充分深入的判断。在评估策略层面，该方法的评分管线对所有输入视频施加统一的评价标准与推理路径，然而不同类型的视频对重要性的界定截然不同，新闻视频侧重信息密度，叙事短片依赖情节推进，生活记录更关注情感表达。固定管线无法在推理之前对视频内容进行全局审视并动态调整评价侧重，仅将大语言模型用作被动的评分函数，未能激发其主动规划与策略决策的潜力。在时序建模层面，该方法以场景为独立单元逐一评分，每个场景的重要性判断均不感知前序场景的具体内容与得分分布。这种无记忆的评分模式在长视频中极易导致两类典型偏差：内容相近的场景被冗余保留，而叙事过渡区间的桥接性场景则因语义密度不高被系统性低估。

上述三重局限的根源在于，第三章方法虽已将大语言模型引入重要性评估环节，但其单模型、固定流程、无记忆的推理架构仅释放了大语言模型作为零样本评分器的基础能力，尚未充分调动其在主动规划、多角色协作与上下文推理方面的深层潜力。

为突破上述瓶颈，本章提出面向视频摘要的多智能体大语言模型推理框架，将重要性评估从单模型统一评分升级为规划、评估和记忆三阶段协同推理流程。具体而言，系统首先将长视频划分为连续片段并生成段级文本描述；随后由 Planner Agent 通览全部片段描述，推断视频主题与叙事结构，并动态分配各评价维度的权重；继而，多个 Expert Agent 分别从叙事推进、视觉显著性、情绪表达与信息密度四个维度对片段进行独立评分，并与 Planner Agent 的全局分数加权融合得到段级得分。在评分过程中，系统引入轻量级记忆机制，将前序片段的内容摘要与评分结论注入后续推理上下文，以增强时序一致性并抑制冗余保留。最终，段级得分映射回帧序列，经归一化与平滑处理后生成帧级摘要结果。

由此可见，本章所讨论的问题并不是单纯的视频内容压缩，而是一个贯通段级语义表征、全局规划、多专家协同推理、上下文记忆建模以及帧级摘要生成的整体优化流程。该流程的目标不是通过额外训练学习统一参数化评分函数，而是借助大语言模型驱动的多角色协作，在开放域视频上实现更具解释性的免训练视频摘要生成。

\xsection{基于大语言模型的多智能体协同的视频摘要方法}{LLM-agent collaborative video summarization method}

针对上一节归纳出的评价维度耦合、评估策略静态以及时序上下文缺失三项问题，本文设计了如图~\ref{fig:c4:method_overview} 所示的多智能体协同视频摘要框架。该框架在实现层面对应一条统一的推理管线：输入为原始视频及其帧级时序信息，输出为与公开数据集评测协议对齐的帧级重要性分数。与第三章单模型逐段打分的推理范式不同，本章方法将视频摘要过程分解为段级表示构建、全局规划与多专家评分、上下文注入与帧级回填三个相互衔接的子过程，从而将片段是否应被保留的判断从单次提示中的隐式折中，转化为具有明确中间状态的结构化推理链。

\begin{figure}[t]
	\centering
	\includegraphics[width=\textwidth]{c4/architecture_c4.pdf}
	\caption{基于大语言模型智能体协同的视频摘要方法总览。}\label{fig:c4:method_overview}
\end{figure}

设输入视频为 $\mathcal{V}= \{f_{1}, f_{2}, \dots, f_{T}\}$，其总帧数为 $T$，数据集给定的采样位置集合为 $\mathcal{Q}=\{q_1,q_2,\dots,q_M\}$。当前实现将图~\ref{fig:c4:method_overview} 所示流程抽象为如下复合映射：

\begin{equation}
	\mathbf{r}= \mathcal{F}_{\mathrm{map}}\Big(\mathcal{F}_{\mathrm{score}}\big(\mathcal{F}_{\mathrm{plan}}(\mathcal{F}_{\mathrm{cap}}(\mathcal{F}_{\mathrm{seg}}(\mathcal{V})))\big),\mathcal{Q},T\Big)
\end{equation}
式中，$\mathbf{r}$——与采样位置对齐的帧级重要性分数序列；$\mathcal{F}_{\mathrm{seg}}$——视频分段映射；$\mathcal{F}_{\mathrm{cap}}$——段级描述生成映射；$\mathcal{F}_{\mathrm{plan}}$——视频级规划映射；$\mathcal{F}_{\mathrm{score}}$——段级评分映射；$\mathcal{F}_{\mathrm{map}}$——段到帧的映射函数；$\mathcal{Q}$——采样位置集合；$T$——视频总帧数。

其中，$\mathcal{F}_{\mathrm{seg}}$ 负责将连续视频帧组织为片段序列，$\mathcal{F}_{\mathrm{cap}}$ 将每个片段压缩为单条文本描述，$\mathcal{F}_{\mathrm{plan}}$ 基于全部片段描述生成视频级规划结果，$\mathcal{F}_{\mathrm{score}}$ 在规划约束下对各片段给出重要性分数，$\mathcal{F}_{\mathrm{map}}$ 则将段级结果确定性地映射为采样帧位置上的帧级分数。由此，整条链路的中间变量可写为

\begin{equation}
	\mathcal{S}=\{S_1,S_2,\dots,S_N\},\quad
	\mathcal{C}=\{(S_1,c_1),(S_2,c_2),\dots,(S_N,c_N)\},\quad
	\hat{\mathbf{y}}=\{\hat{y}_1,\hat{y}_2,\dots,\hat{y}_N\}
\end{equation}
式中，$\mathcal{S}$——段级结构集合；$S_k$——第 $k$ 个片段；$\mathcal{C}$——段级语义描述集合；$c_k$——第 $k$ 个片段的文本描述；$\hat{\mathbf{y}}$——段级重要性分数序列；$N$——片段总数。

其中 $\mathcal{S}$、$\mathcal{C}$ 与 $\hat{\mathbf{y}}$ 分别对应段级结构、段级语义描述以及段级重要性结果。当前实现的核心特征不在于引入额外训练，而在于以显式的数据流定义支撑后续各模块的职责划分：前端负责降低视觉冗余，中端负责构造多维语义证据，后端负责将评分结果回填到评测所需的时间尺度。

基于上述定义，4.2 节围绕三个紧密耦合的模块展开：4.2.1 给出视频结构化处理与段级描述生成机制，用以构造可被语言模型消费的输入接口；4.2.2 给出 Planner-Expert 协同评分机制，用以将全局语境与局部证据结合为段级分数；4.2.3 给出轻度上下文注入与帧级映射机制，用以把段级推理结果转换为最终摘要输出。由于这些模块在代码中由统一管线顺序驱动，其关系并非概念上的并列堆叠，而是前一模块的输出即后一模块的输入。


\xsubsection{视频结构化处理与段级描述生成}{Video Structuring and Segment-level Caption Generation}

该阶段的功能是将高维帧流压缩为有序的段级文本序列，并为后续规划与评分提供统一输入接口。由于当前实现中的 Planner 与各 Expert 均以文本提示为输入，若直接将原始视频帧送入评分模块，不仅会导致上下文长度迅速膨胀，还会使大量弱信息帧在推理中占据不必要的表示预算。基于这一约束，系统并不直接在帧级空间进行重要性判断，而是先在时间轴上构造片段，再为每个片段生成单条语义描述，使视频内容由像素序列转化为可排序、可追踪的文本单元序列。

在实现上，视频读入模块首先返回视频标识、帧率与总帧数等基础信息，并据此构造片段集合 $\mathcal{S}= \{S_{1}, S_{2}, \dots, S_{N}\}$。第 $k$ 个片段定义为

\begin{equation}
	S_{k}= \{k,b_{k},e_{k},\mathcal{P}_{k}\}
\end{equation}
式中，$S_k$——第 $k$ 个片段的结构化表示；$k$——片段编号；$b_k$、$e_k$——片段在原视频中的起止帧边界；$\mathcal{P}_k$——片段内部的采样帧索引集合。

其中 $k$ 为片段编号，$b_k$ 与 $e_k$ 分别表示片段在原视频中的起止帧边界，$\mathcal{P}_{k}=\{p_{k,1},p_{k,2},\dots,p_{k,m_k}\}$ 为该片段内部用于描述生成的采样帧索引。当前实现支持三种分段模式：按固定段数划分、按固定窗口长度划分以及带重叠的滑动窗口划分。若采用按段数划分，则片段长度近似为

\begin{equation}
	\ell = \max\left(1, \left\lfloor \frac{T}{N} \right\rfloor\right)
\end{equation}
式中，$\ell$——按段数切分时的近似片段长度；$T$——视频总帧数；$N$——预设片段数量。

若采用固定窗口或滑动窗口模式，设窗口长度为 $L$、重叠长度为 $o$，则相邻片段的推进步长为

\begin{equation}
	\Delta = L - o
\end{equation}
式中，$\Delta$——相邻片段的推进步长；$L$——窗口长度；$o$——相邻窗口之间的重叠长度。

上述定义表明，当前实现中的 segment 是服务于推理组织的时间单元，而非严格意义上的镜头或事件边界。其设计目标并非恢复最细粒度的叙事结构，而是在不引入复杂分割学习模块的前提下，为后续多智能体推理提供长度可控、顺序明确且覆盖全视频的候选单元。

在完成时间结构化之后，系统对每个片段执行均匀采样，而非对全部帧逐一描述。对长度为 $|S_k|=e_k-b_k$ 的片段，若预设采样数为 $m_k$，则采样位置写为

\begin{equation}
	p_{k,j}= \operatorname{round}\left(b_k + \frac{(j-1)(|S_k|-1)}{m_k-1}\right), \quad j=1,2,\dots,m_k
\end{equation}
式中，$p_{k,j}$——第 $k$ 个片段内第 $j$ 个采样帧的位置；$|S_k|$——片段 $S_k$ 的长度；$m_k$——该片段的预设采样帧数；$\operatorname{round}(\cdot)$——四舍五入操作。

当片段长度不足以支撑预设采样数时，实际采样数会被截断到片段长度；当仅保留一个采样位置时，系统退化为单帧采样。均匀采样并不试图恢复片段内部的全部时序细节，而是以尽可能低的视觉输入成本覆盖片段的主要变化趋势，从而使后续视觉语言模型获得的是该片段的整体概貌，而非某一时刻的局部快照。

\begin{figure}[t]
	\centering
	\includegraphics[width=0.7\textwidth]{c4/model_c4_01.pdf}
	\caption{视频结构化处理与段级描述生成流程示意图。}\label{fig:c4:segment_caption_pipeline}
\end{figure}

采样帧提取完成后，系统将同一片段内的多帧图像联合输入视觉语言模型，并为该片段生成单条文本描述。若记片段 $S_k$ 对应的图像集合为 $I_k=\{I_{k,1},I_{k,2},\dots,I_{k,m_k}\}$，则段级描述生成过程可写为

\begin{equation}
	c_k = \Phi_{\mathrm{VLM}}\big(I_k; \theta_{\mathrm{vlm}}, \pi_{\mathrm{seg}}\big)
\end{equation}
式中，$c_k$——第 $k$ 个片段的段级描述；$\Phi_{\mathrm{VLM}}(\cdot)$——视觉语言模型描述函数；$I_k$——片段 $S_k$ 对应的图像集合；$\theta_{\mathrm{vlm}}$——视觉语言模型参数；$\pi_{\mathrm{seg}}$——段级提示模板。

其中 $\theta_{\mathrm{vlm}}$ 表示视觉语言模型参数，$\pi_{\mathrm{seg}}$ 表示段级描述提示模板。当前实现使用冻结的本地视觉语言模型对多帧联合推理，并将输出限制为单段短文本，其关注点被约束为持续动作、场景上下文以及主要事件。与第三章的逐帧描述再聚合不同，这里直接生成段级描述 $c_k$，其结果与片段边界、采样帧索引共同构成结构化对象

\begin{equation}
	\tilde{S}_{k}=\{k,b_k,e_k,\mathcal{P}_k,c_k\}
\end{equation}
式中，$\tilde{S}_k$——同时包含边界、采样信息和文本描述的增强片段对象；$c_k$——第 $k$ 个片段的段级文本描述。

这一设计在实现上的直接作用，是使后续规划与评分模块面向统一的文本接口工作，而不再依赖多级视觉特征聚合。

需要指出的是，段级描述生成阶段还包含一项面向可用性的退化处理：当视频读取失败、采样为空或模型推理异常时，系统返回通用兜底文本，以保证输出序列在结构上完整。该设计不改变方法主体的语义能力，却能保证后续规划与评分过程始终作用于长度一致的片段序列。

综上，前端模块的输出并非抽象的视频理解结果，而是带有片段边界与采样信息的段级语义序列

\begin{equation}
	\mathcal{C}= \{(S_{1}, c_{1}), (S_{2}, c_{2}), \dots, (S_{N}, c_{N})\}
\end{equation}
式中，$\mathcal{C}$——段级语义序列；$S_k$——第 $k$ 个片段；$c_k$——与该片段对应的文本描述；$N$——片段总数。

其中每个元素同时保留片段边界、采样帧索引与文本描述。该表示既维持了时间顺序，又显著压缩了视觉输入规模，从而为下一阶段的全局规划与多专家评分提供可直接消费的输入基础。

\xsubsection{Planner-Expert 协同评分与分数生成机制}{Planner-Expert Collaborative Scoring Mechanism}

第二阶段的目标是将段级描述序列转化为可比较的段级重要性分数。当前实现采用两步式评分结构：一方面，Planner Agent 基于全视频描述建立全局语境；另一方面，Planner 与四类 Expert 在统一语境约束下对每个片段进行独立评分，并通过显式加权生成最终段级得分。该设计的关键不在于增加评分器数量，而在于将全局任务定义与局部证据判断解耦，以避免单次提示同时承担主题识别、摘要概括和局部打分三种职责。

在全局规划步骤中，Planner Agent 接收全部片段描述 $\{c_1,c_2,\dots,c_N\}$，输出视频级规划对象

\begin{equation}
	\mathcal{P}= \{z, g, \mathbf{w}, \rho\}
\end{equation}
式中，$\mathcal{P}$——Planner 输出的视频级规划对象；$z$——视频主题标签；$g$——全局摘要；$\mathbf{w}$——专家权重向量；$\rho$——Planner 给出的文本化规划依据。

其中 $z$ 表示视频主题标签，$g$ 表示面向整段视频的全局摘要，$\mathbf{w}=[w_s,w_v,w_e,w_i]$ 表示四类 Expert 的动态权重，$\rho$ 为 Planner 返回的文本化规划依据。实现中，$z$ 从有限主题集合中选取，用于稳定后续提示中的任务语义；$g$ 为片段评分提供跨片段语境；$\mathbf{w}$ 则决定不同评分维度在融合时的相对贡献。为保证不同维度可以在同一尺度下加权，权重经归一化后满足

\begin{equation}
	\sum_{m \in \{s,v,e,i\}} w_m = 1, \quad w_m \ge 0
\end{equation}
式中，$w_m$——第 $m$ 类专家的融合权重；$m\in\{s,v,e,i\}$——分别对应 story、visual、emotion 和 information 四类专家。

其中 $m\in\{s,v,e,i\}$ 分别对应 story、visual、emotion 与 information 四个维度。当模型未返回有效权重时，系统退化为均匀分配，以保持后续融合过程的数值可控性。

在获得 $\mathcal{P}$ 之后，系统对每个片段执行逐段评分。对当前片段 $S_k$ 及其描述 $c_k$，Planner Agent 基于 $z$、$g$、$c_k$ 以及历史上下文生成一个主分数 $p_k$；四个 Expert Agent 则分别从不同语义维度给出补充分数 $e_k^{(s)}$、$e_k^{(v)}$、$e_k^{(e)}$ 与 $e_k^{(i)}$。各 Expert 的职责划分如下：

\begin{itemize}
	\item \textbf{Story Expert}：衡量当前片段对叙事推进或核心事件演化的贡献程度。
	\item \textbf{Visual Expert}：衡量当前片段在视觉上是否具有高光性、显著性或场景变化价值。
	\item \textbf{Emotion Expert}：衡量当前片段中情绪表达、张力变化或情感感染力的强弱。
	\item \textbf{Information Expert}：衡量当前片段提供的新信息、讲解价值或语义增量。
\end{itemize}

需要强调的是，多 Expert 结构在当前实现中并不对应多模型交互或多轮协商，而是同一类语言模型接口在不同角色化提示下的并行调用。因而，这一模块的贡献不在于形成复杂的 Agent 对话，而在于将局部片段的重要性依据拆解为叙事、视觉、情绪与信息四类可解释证据。

在数值层面，Planner 与各 Expert 返回的基础分数均被裁剪到 $[0,1]$ 区间。设裁剪函数为 $\psi(\cdot)$，则有

\begin{equation}
	p_k = \psi(\tilde{p}_k), \quad e_k^{(m)} = \psi(\tilde{e}_k^{(m)}), \quad \psi(x)= \min\{1, \max(0,x)\}
\end{equation}
式中，$p_k$——第 $k$ 个片段的 Planner 主分数；$e_k^{(m)}$——第 $m$ 类 Expert 对第 $k$ 个片段给出的分数；$\tilde{p}_k$、$\tilde{e}_k^{(m)}$——模型原始输出；$\psi(\cdot)$——分数裁剪函数。

其中 $\tilde{p}_k$ 与 $\tilde{e}_k^{(m)}$ 为模型原始输出。基础分数裁剪与权重归一化共同构成了当前实现中最核心的数值约束：前者抑制异常输出，后者确保多维证据的组合具有稳定尺度。

在此基础上，段级最终分数定义为

\begin{equation}
	\hat{y}_{k}= p_{k}+ \sum_{m \in \{s,v,e,i\}} w_{m} e_{k}^{(m)}
\end{equation}
式中，$\hat{y}_k$——第 $k$ 个片段的最终重要性分数；$p_k$——Planner 主分数；$w_m$——第 $m$ 类专家的权重；$e_k^{(m)}$——对应专家分数。

该定义表明，Planner 输出并非与 Expert 完全对等的附加分量，而是承担全局主导项的角色；四类 Expert 则提供受主题与摘要约束的局部补充证据。由于 $p_k$ 与 $e_k^{(m)}$ 均位于 $[0,1]$ 区间，而专家权重和为 1，最终分数 $\hat{y}_k$ 不再被强制约束到 $[0,1]$，而是保留相对尺度信息，供后续帧级映射与标准化处理使用。这一处理与代码实现保持一致，也使得段间排序在融合后不必再次经历额外裁剪。

\begin{figure}[t]
	\centering
	\includegraphics[width=0.8\textwidth]{c4/model_c4_02.pdf}
	\caption{Planner-Expert 协同评分与分数生成机制示意图。}\label{fig:c4:planner_expert_scoring}
\end{figure}

在实现结果上，每个片段最终形成一条结构化评分记录，其中除 $\hat{y}_k$ 外，还保留 Planner 主分数、四类 Expert 分数及其对应理由。因而，第二阶段的输出并不是简单的标量序列，而是兼具数值结果与语义证据的段级评分对象。对于论文表述而言，这一点意味着本章方法的可解释性并非来自事后分析，而是内嵌于评分数据结构本身。



\xsubsection{轻度上下文记忆注入机制}{Lightweight Contextual Memory Injection}

前两小节解决了如何构造片段表示与如何形成段级分数两个问题，但最终摘要输出仍需回答两个实现层面的关键问题：其一，局部评分如何感知前序语境；其二，段级分数如何映射到公开数据集要求的帧级表示。当前实现将这两个问题统一收束在第三阶段完成，即在逐段评分过程中注入轻度上下文记忆，并在评分结束后执行确定性的帧级回填与评测对齐。

所谓轻度记忆，并非额外训练出的状态更新网络，而是对历史片段描述的有限窗口拼接。设当前片段为 $S_k$，其可访问的历史描述集合记为

\begin{equation}
	\mathcal{M}_k = \{c_{1}, c_{2}, \dots, c_{k-1}\}
\end{equation}
式中，$\mathcal{M}_k$——当前片段 $S_k$ 可访问的全部历史描述集合；$c_1,\dots,c_{k-1}$——当前片段之前已经生成的段级描述。

当最大历史窗口长度设为 $h$ 时，实际送入提示的上下文为

\begin{equation}
	\mathcal{M}_k^{(h)} = \{c_{k-h}, c_{k-h+1}, \dots, c_{k-1}\}
\end{equation}
式中，$\mathcal{M}_k^{(h)}$——长度受限的历史上下文集合；$h$——最大历史窗口长度。

其中越界项自动忽略；若记忆机制关闭，则 $\mathcal{M}_k^{(h)}$ 为空。实现上，系统将这些历史描述按时间顺序线性拼接为文本上下文，并在 Planner 与各 Expert 评分当前片段时一并输入。因此，该机制建模的并不是抽象的长期依赖，而是最近若干已处理片段对当前判断的约束。这一设计的直接作用是抑制相邻重复内容被连续高估，并为叙事过渡片段提供必要的前文语义参照。

在得到全部段级分数 $\hat{\mathbf{y}}=\{\hat{y}_1,\hat{y}_2,\dots,\hat{y}_N\}$ 后，系统需要将这些结果映射回帧轴。当前实现采用确定性的区间回填策略：每个片段的最终分数被赋值给其覆盖区间内的全部帧；若某一帧同时落在多个片段内，则取所有覆盖片段分数的平均值。令 $\Omega_t$ 表示覆盖第 $t$ 帧的片段集合，则稠密帧分数定义为

\begin{equation}
	s_t = \frac{1}{|\Omega_t|}\sum_{k \in \Omega_t} \hat{y}_k
\end{equation}
式中，$s_t$——第 $t$ 帧的稠密分数；$\Omega_t$——覆盖第 $t$ 帧的片段集合；$|\Omega_t|$——集合 $\Omega_t$ 中的片段数量；$\hat{y}_k$——第 $k$ 个片段的最终重要性分数。

对于无重叠分段，上式退化为单片段直接赋值；对于带重叠的滑动窗口模式，则该定义自然完成对重叠区域的平均融合。因此，第三阶段中的帧级映射并不是额外的学习模块，而是由片段边界和段级分数共同确定的后处理过程。

由于公开数据集通常仅在采样位置上提供评分与评测元数据，系统进一步依据采样位置集合 $\mathcal{Q}=\{q_1,q_2,\dots,q_M\}$ 提取对齐后的帧级分数

\begin{equation}
	r_i = s_{q_i}, \quad i=1,2,\dots,M
\end{equation}
式中，$r_i$——第 $i$ 个采样位置对应的帧级分数；$s_{q_i}$——原始帧轴上位置 $q_i$ 的稠密分数；$M$——采样位置总数。

由此得到输出序列 $\mathbf{r}=\{r_1,r_2,\dots,r_M\}$。这一序列正是整条推理管线在方法层面的最终输出，也是后续评测模块的直接输入。

需要进一步说明的是，当前实现中的 $\mathbf{r}$ 在进入数据集评测前还会经历标准化与可选平滑。记这两种评测变体分别为 $\mathbf{r}^{\mathrm{raw}}$ 与 $\mathbf{r}^{\mathrm{smooth}}$，则有

\begin{equation}
	\mathbf{r}^{\mathrm{raw}} = \operatorname{Norm}(\mathbf{r}),
	\quad
	\mathbf{r}^{\mathrm{smooth}} = \operatorname{Smooth}(\operatorname{Norm}(\mathbf{r});w)
\end{equation}
式中，$\mathbf{r}^{\mathrm{raw}}$——归一化后的原始帧级分数序列；$\mathbf{r}^{\mathrm{smooth}}$——平滑后的帧级分数序列；$\operatorname{Norm}(\cdot)$——区间归一化操作；$\operatorname{Smooth}(\cdot)$——平滑操作；$w$——平滑窗口长度。

其中 $\operatorname{Norm}(\cdot)$ 为区间归一化，$\operatorname{Smooth}(\cdot)$ 为边界保持的滑动平均，$w$ 为平滑窗口长度。由于上一阶段得到的 $\hat{y}_k$ 保留了相对尺度而未作额外归一化，这一步承担了将不同视频、不同分段设置的分数拉回统一评测空间的作用。

在公开数据集协议下，仅有帧级分数仍不足以生成最终摘要，因为标准评测要求满足长度预算的二值摘要向量。为此，系统依据 change points 形成候选评测分段，并在总长度约束下求解 0/1 背包优化。若第 $j$ 个候选分段的平均分数为 $u_j$、长度为 $l_j$，容量上限为 $B$，则摘要选择可表示为

\begin{equation}
	\max_{\mathbf{z}} \sum_{j=1}^{K} z_j u_j \quad
	\text{s.t.} \quad \sum_{j=1}^{K} z_j l_j \le B, \quad z_j \in \{0,1\}
\end{equation}
式中，$\mathbf{z}$——候选评测分段的选择向量；$z_j$——第 $j$ 个候选分段是否被选入摘要的决策变量；$u_j$——第 $j$ 个候选分段的平均分数；$l_j$——第 $j$ 个候选分段的长度；$B$——摘要长度预算；$K$——候选评测分段总数。

其中 $K$ 表示候选评测分段数，$B$ 由视频总帧数与摘要比例共同确定。至此，第三阶段完成了从段级语义判断到帧级评测输入、再到二值摘要构造的完整对接。

综上，4.2.3 的核心并非单独引入一个独立模块，而是在实现层面把历史文本上下文、段级分数回填、采样对齐以及标准摘要构造组织为统一的后端过程。由此，4.2.1 生成的段级描述和 4.2.2 生成的段级分数才真正转化为可与 SumMe、TVSum 等公开数据集协议直接对接的摘要输出。

至此，图~\ref{fig:c4:method_overview} 所示方法链路已完整闭合：前端通过结构化分段与段级描述生成完成输入压缩，中端通过 Planner-Expert 机制完成段级重要性推理，后端通过轻度上下文注入与帧级映射完成结果对齐。三者共同构成当前代码实现下的多智能体视频摘要核心方法。下一节将在此基础上对实验设置与性能表现进行分析。

\xsection{实验与分析}{Experiments and Analysis}

\xsubsection{数据集与评估协议}{Datasets and Evaluation Protocol}

本节的目的在于说明第四章方法所采用的实验对象、数据组织形式以及与当前工程实现严格一致的评测协议。由于本章方法最终输出的是可与标准视频摘要基准直接对接的帧级分数序列，因此实验设置必须同时满足两个要求：其一，数据集需要提供视频、采样位置、分段边界以及人工摘要或人工评分等基础信息；其二，评测流程必须能够将本章方法生成的段级结果稳定映射到公开基准所要求的帧级评价形式。依据当前实现，系统支持的基准数据集为 SumMe 与 TVSum，并围绕它们构建统一的数据读取与评测接口。

从数据组织角度看，当前实现并不直接从松散标注文件中构建评测样本，而是通过预先整理好的 HDF5 数据文件与名称映射文件加载单视频记录。对于给定视频，数据读取模块会解析视频路径、总帧数、采样位置、change points、每段帧数、人类摘要以及人类评分等信息，并将其封装为统一的数据记录对象。若记某个视频样本为 $\mathcal{D}_v$，则在当前实现中可抽象表示为：

\begin{equation}
	\mathcal{D}_v = \{V, T, \mathcal{Q}, \mathcal{B}, \mathcal{L}, \mathcal{Y}_{\mathrm{human}}, \mathcal{R}_{\mathrm{human}}\}
\end{equation}
式中，$\mathcal{D}_v$——视频样本 $v$ 的统一数据记录对象；$V$——视频文件路径；$T$——视频总帧数；$\mathcal{Q}$——采样帧位置集合；$\mathcal{B}$——change points 定义的评测分段边界；$\mathcal{L}$——各评测分段长度；$\mathcal{Y}_{\mathrm{human}}$——人类摘要；$\mathcal{R}_{\mathrm{human}}$——人类评分。

其中 $V$ 表示视频文件路径，$T$ 表示视频总帧数，$\mathcal{Q}$ 表示采样帧位置集合，$\mathcal{B}$ 表示 change points 所定义的评测分段边界，$\mathcal{L}$ 表示各评测分段的长度，$\mathcal{Y}_{\mathrm{human}}$ 表示人类摘要，$\mathcal{R}_{\mathrm{human}}$ 表示人类评分。该统一数据结构使得后续推理与评测流程不再依赖数据集内部文件格式差异，而能够在统一接口上完成实验管理。

当前实验对象包括两个公开视频摘要基准数据集。第一类是 SumMe，系统通过对应的 HDF5 文件与映射文件读取其视频及人工摘要信息；第二类是 TVSum，系统以相同方式读取视频路径、采样点、分段信息以及人工评分或人工摘要。需要强调的是，本章在这一小节中只陈述当前代码可直接确认的数据接口与支持范围，而不额外引入未在实现中显式使用的统计描述。因此，这里不再展开诸如视频类别数量、平均时长等扩展性数据集介绍；若后续需要在论文中补充更详细的数据集背景，应以官方数据集文档或原始论文为准。

\begin{figure}[t]
	\centering
	\fbox{\parbox[c][4.8cm][c]{0.92\linewidth}{\centering
		图片占位：数据集样本结构示意图。\\
		建议展示单个视频样本在当前代码中的字段组织形式，包括 video\_path、n\_frames、picks、change\_points、n\_frame\_per\_seg、user\_summary 与 user\_scores 等内容，以及这些字段与后续评测流程的对应关系。
	}}
	\caption{实验数据记录结构与评测字段示意图。}\label{fig:c4:dataset_record_structure}
\end{figure}

在评测协议方面，当前实现沿用了视频摘要任务中先构造机器摘要、再与人工摘要或人工评分比较的标准流程。给定本章方法输出的帧级分数序列 $\mathbf{r}= \{r_1, r_2, \dots, r_M\}$，系统首先依据采样位置将其扩展回完整帧轴，并结合数据集提供的 change points 与各分段长度，构造满足长度约束的机器摘要向量 $\mathbf{y}_{\mathrm{sys}}$。随后，将该机器摘要与人工摘要 $\mathbf{y}_{\mathrm{gt}}$ 比较，得到基于重叠度的评价指标。若记精确率与召回率分别为 $P$ 和 $R$，则当前实现中的核心重叠性指标定义为：

\begin{equation}
	P = \frac{|\mathbf{y}_{\mathrm{sys}} \cap \mathbf{y}_{\mathrm{gt}}|}{|\mathbf{y}_{\mathrm{sys}}| + \varepsilon},
	\quad
	R = \frac{|\mathbf{y}_{\mathrm{sys}} \cap \mathbf{y}_{\mathrm{gt}}|}{|\mathbf{y}_{\mathrm{gt}}| + \varepsilon}
\end{equation}
式中，$P$——精确率；$R$——召回率；$\mathbf{y}_{\mathrm{sys}}$——机器摘要向量；$\mathbf{y}_{\mathrm{gt}}$——人工摘要向量；$\varepsilon$——数值稳定项。

\begin{equation}
	F_{1}= \frac{2PR}{P+R}
\end{equation}
式中，$F_1$——精确率与召回率的调和平均指标；$P$——精确率；$R$——召回率。

其中 $\varepsilon$ 为数值稳定项。除基于二值摘要的重叠指标外，当前实现还基于帧级预测分数与人类评分序列之间的一致性，计算 Spearman 秩相关系数 $\rho$ 与 Kendall 秩相关系数 $\tau$，用于衡量模型对相对重要性排序的保持能力。

对于多用户标注情形，当前代码对 SumMe 与 TVSum 采用了不同的聚合策略。对于 SumMe，系统使用最大 $F_1$ 聚合方式；对于 TVSum，则采用平均 $F_1$ 聚合方式。形式化地，若同一视频对应 $U$ 个人工摘要，则其最终 $F_1$ 可写为：

\begin{equation}
	F_{1}^{\mathrm{SumMe}} = \max_{u \in \{1,\dots,U\}} F_{1}^{(u)},
	\quad
	F_{1}^{\mathrm{TVSum}} = \frac{1}{U}\sum_{u=1}^{U} F_{1}^{(u)}
\end{equation}
式中，$F_{1}^{\mathrm{SumMe}}$——SumMe 数据集上的最终 $F_1$ 聚合结果；$F_{1}^{\mathrm{TVSum}}$——TVSum 数据集上的最终 $F_1$ 聚合结果；$F_{1}^{(u)}$——相对于第 $u$ 位标注者计算得到的 $F_1$ 分数；$U$——标注者数量。

这一设置与当前实现中的 reduction 策略保持一致，并体现了两个数据集在标准评测协议上的差异。

除了基础评测逻辑外，当前实验接口还在正式计算指标之前，对模型输出帧级分数构造两个评测变体，即归一化后的原始分数序列与归一化后进一步平滑的分数序列。记二者分别为 $\mathbf{r}^{\mathrm{raw}}$ 与 $\mathbf{r}^{\mathrm{smooth}}$，则当前评测流程会对这两个变体分别计算全部指标。这一设计的目的不是引入新的方法模块，而是便于在实验阶段区分原始推理结果与简单后处理策略对摘要质量的影响。由于当前工作区中尚未提供真实实验输出文件，本小节仅说明这一协议设置的存在，不对其优劣作任何定量结论。

\begin{figure}[t]
	\centering
	\fbox{\parbox[c][5.0cm][c]{0.92\linewidth}{\centering
		图片占位：评测协议流程图。\\
		建议展示从 frame scores 输入开始，经归一化、平滑、change points 分段、knapsack 摘要构造，到 F1、Precision、Recall、Spearman 和 Kendall 计算的整体评测流程。
	}}
	\caption{当前实现中的视频摘要评测协议流程示意图。}\label{fig:c4:evaluation_protocol}
\end{figure}

需要说明的是，当前实现还支持两种实验运行方式：一种是针对单个指定视频执行推理与评测；另一种是针对整个数据集逐视频运行，并基于官方 split 信息进一步做聚合统计。这意味着第四章后续实验既可以汇报单视频案例分析，也可以汇报按数据集或按 split 聚合后的整体表现。然而，由于当前工作区中未发现现成的实验输出文件、汇总报告或可视化结果，本小节不对具体数值、表格和性能结论进行描述，相关内容将在后续实际运行实验后补充。

综上，本小节明确了第四章实验部分所依赖的两个基准数据集与当前代码实现中的评测协议：系统以 SumMe 和 TVSum 作为标准测试对象，通过统一的数据记录接口读取视频与标注信息，通过帧级分数映射和机器摘要构造实现与官方协议对齐，并最终采用 $F_1$、Precision、Recall、Spearman $\rho$ 和 Kendall $\tau$ 等指标对方法进行量化评估。下一小节将在这一协议基础上，进一步介绍实验所采用的运行设置与关键实现细节。

\xsubsection{实验设置}{Experimental Settings}

在明确数据集与评测协议之后，本小节进一步说明第四章实验所采用的运行方式、模型配置、关键超参数以及输出组织形式。需要强调的是，以下内容均严格依据当前工程主执行流程与配置接口进行整理，其作用是说明实验如何被运行，而非提前给出任何性能结论。由于当前工作区中尚未提供完整的硬件环境记录与既有实验日志，本文仅陈述代码中可以直接确认的设置项，对于无法从实现中直接验证的环境细节不作扩展性描述。

从运行入口看，当前系统通过统一命令行接口完成单视频推理、单视频评测以及整合式运行评测三类实验任务。若记实验任务类型为 $a$，则当前实现支持的任务集合为：

\begin{equation}
	\mathcal{A}= \{\mathrm{run}, \mathrm{eval}, \mathrm{run\_eval}\}
\end{equation}
式中，$\mathcal{A}$——系统支持的实验任务集合；$\mathrm{run}$——仅执行推理；$\mathrm{eval}$——仅执行评测；$\mathrm{run\_eval}$——串联执行推理与评测。

其中，$\mathrm{run}$ 表示仅对输入视频执行分段、描述生成、多智能体评分与帧级分数输出；$\mathrm{eval}$ 表示在已有中间结果基础上执行标准评测；$\mathrm{run\_eval}$ 则表示将推理与评测串联为一次完整实验流程。对于公开数据集上的批量实验，系统还支持指定数据集名称、数据根目录、split 文件目录以及 split 数量，从而逐视频完成推理与评测，并在数据集层面进一步生成聚合结果。由此可见，当前实验框架既可服务于单视频案例分析，也可用于标准基准下的批量统计评估。

在大语言模型调用方式上，当前实现提供真实接口调用与调试用模拟接口两种模式。若记语言模型工作模式为 $m$，则有：

\begin{equation}
	m \in \{\mathrm{api}, \mathrm{mock}\}
\end{equation}
式中，$m$——大语言模型的工作模式；$\mathrm{api}$——真实接口调用模式；$\mathrm{mock}$——调试用模拟接口模式。

其中，$\mathrm{api}$ 模式通过 OpenAI 兼容接口访问外部大语言模型服务，运行时需要提供环境变量中的 API Key，并可选指定自定义 Base URL；$\mathrm{mock}$ 模式则返回结构化的模拟结果，主要用于流程调试、接口连通性验证与单元测试场景。在当前默认设置下，Planner 与 Expert Agent 所共享的语言模型名称为 gpt-4o-mini，温度参数设置为 $0.2$，并保留有限次重试机制以提升接口调用稳定性。需要指出的是，这一设置主要作用于规划与评分阶段，而前端段级描述生成并不依赖该文本接口。

与 Planner 和 Expert 的文本推理模块相对应，段级多模态描述生成阶段采用本地视觉语言模型完成。当前实现默认使用 Qwen3-VL-8B-Instruct 作为 segment captioner 的底层模型，并通过自动设备选择机制加载模型参数。对于每个片段，系统根据预设提示词引导模型关注持续动作、场景上下文与主要事件，并将最大生成长度限制为 96 个新 token。该设置反映出当前实验框架对前端语义压缩模块的基本取舍：一方面保留多帧联合理解能力，另一方面通过较短输出控制描述冗长度，以避免后续 Planner 输入上下文过度膨胀。

在视频结构化参数方面，当前代码支持按段数切分、固定窗口切分和滑动窗口切分三种分段模式。若将分段方式记为 $\mu$，则有：

\begin{equation}
	\mu \in \{\mathrm{count}, \mathrm{fixed\_frames}, \mathrm{sliding\_window}\}
\end{equation}
式中，$\mu$——视频分段方式；$\mathrm{count}$——按段数切分模式；$\mathrm{fixed\_frames}$——按固定帧窗口切分模式；$\mathrm{sliding\_window}$——滑动窗口切分模式。

其中，当前 README 所给出的推荐设置采用按段数切分模式，即默认将每个视频划分为 8 个片段；与此同时，每个片段默认抽取 4 帧用于段级描述生成。若需要启用滑动窗口机制，系统还允许进一步指定窗口重叠长度。上述设置说明，第四章实验中的前端时间结构化并不是固定不变的，而是通过统一参数接口控制；不过在后续实际实验中，若无特别说明，宜优先沿用当前主流程推荐的 count 模式与 8 段设置，以保证全文实验口径一致。

除了分段与描述生成参数外，当前实验还可选择是否启用历史记忆机制。若记记忆开关为 $\delta \in \{0,1\}$，则当 $\delta = 1$ 时，系统会在当前片段评分前拼接前序片段描述作为文本上下文，并通过最大历史片段数 $h$ 控制记忆窗口长度；在 README 给出的推荐运行命令中，$h$ 默认取 8。该设置意味着，当前实验中的记忆机制不是无限累积的长程缓存，而是受窗口约束的轻量上下文注入策略。对于长视频而言，这种截断式设计有利于在保持时序连续性的同时避免提示词长度随视频长度线性增长。

在评测后处理参数方面，当前实现会针对模型输出的帧级分数生成两个评测变体，即归一化原始分数与归一化平滑分数。设原始帧级分数序列为 $\mathbf{r}$，则可进一步得到：

\begin{equation}
	\mathbf{r}^{\mathrm{raw}} = \operatorname{Norm}(\mathbf{r}),
	\quad
	\mathbf{r}^{\mathrm{smooth}} = \operatorname{Smooth}(\operatorname{Norm}(\mathbf{r}); w)
\end{equation}
式中，$\mathbf{r}^{\mathrm{raw}}$——评测使用的归一化原始分数序列；$\mathbf{r}^{\mathrm{smooth}}$——评测使用的归一化平滑分数序列；$w$——平滑窗口长度。

其中 $w$ 表示平滑窗口长度。当前命令行接口中，该窗口默认取 5。实验过程中，系统会分别对这两个变体计算 $F_1$、Precision、Recall、Spearman $\rho$ 与 Kendall $\tau$，以观察简单后处理对最终摘要质量的影响。需要说明的是，这一平滑操作属于评测前的信号整形步骤，而非第四章方法主体中新增的学习模块，因此后文在讨论实验结果时应将其视为评测设置的一部分。

对于整数据集实验，当前实现还支持基于 split 的聚合统计。具体而言，系统可从给定的 split 根目录中读取预定义划分，并按照设定的 split 数量逐一汇总不同 split 上的实验结果。在当前命令行接口中，split 数量被约束为 5 或 50 两种常见配置。每次聚合完成后，系统会进一步生成面向整体实验的统计文件，包括 JSON 格式的指标概览、逐视频记录以及 Markdown 形式的汇总表。该设计有助于将逐视频推理结果自动整理为可直接分析的实验档案，从而减少人工汇总误差。

\begin{figure}[t]
	\centering
	\fbox{\parbox[c][4.8cm][c]{0.92\linewidth}{\centering
		图片占位：实验运行配置示意图。\\
		建议展示命令行入口、任务类型、LLM 模式、caption 模型、分段参数、记忆开关和平滑窗口等配置项之间的关系，并标注单视频运行与数据集批量运行两种路径。
	}}
	\caption{当前实现中的实验运行配置与参数组织示意图。}\label{fig:c4:experiment_settings}
\end{figure}

从实验输出组织形式看，当前系统将推理结果与评测结果分别写入输出目录与评测目录。在单视频评测场景下，系统会按实验名、数据集名和视频名的层次创建子目录，并保存帧级分数、不同评测变体结果、预测与人工曲线对比图以及整体概览文件。进一步地，系统还会自动维护 overview.json、overview\_records.json 与 overview.md 等汇总文件，以便后续在论文中整理表格与案例分析。因此，第四章实验部分不仅具备标准化的运行接口，也已经具备较完整的结果落盘与汇总能力。

\begin{figure}[t]
	\centering
	\fbox{\parbox[c][4.8cm][c]{0.92\linewidth}{\centering
		图片占位：实验输出目录与报告文件示意图。\\
		建议展示 outputs 与 outputs/evaluation 下的目录层次，以及 frame\_scores\_variants.json、eval\_normalized\_raw.json、eval\_normalized\_smoothed.json、overview.json 和 overview.md 等文件之间的关系。
	}}
	\caption{实验结果落盘与汇总报告组织形式示意图。}\label{fig:c4:experiment_outputs}
\end{figure}

综上，第四章实验采用统一命令行入口组织推理与评测流程，在模型侧结合 gpt-4o-mini 与 Qwen3-VL-8B-Instruct 分别完成文本推理和段级描述生成，在参数侧提供分段模式、采样帧数、记忆窗口与平滑窗口等可控配置，并在结果侧形成按视频、按数据集和按 split 的多层次输出结构。这些设置共同构成了后续实验分析的实现基础。下一小节将在此基础上，对方法的整体性能表现进行进一步讨论。

\xsubsection{总体性能表现}{Overall Performance}

本小节的目的在于给出第四章方法在公开基准上的总体性能表现，并以规范化表格形式预留与已有代表性方法的对比位置。依据前一小节给出的实验协议，本文优先采用按官方 split 聚合后的 Mean $F_1$ 作为主指标，同时在必要时结合 Precision、Recall 以及秩相关指标辅助分析。需要说明的是，当前工作区中已经可以直接核验本方法的多组真实实验结果，但尚未提供统一整理完成的文献基线数值表。因此，表\ref{tab:c4:overall_performance} 中仅填写当前仓库内已落盘且可核验的本方法结果，其余方法保留为待补位置，后续应结合正式引用文献逐项补全。

为与第三章实验写法保持一致，表\ref{tab:c4:overall_performance} 按方法类别组织对比对象，并统一汇报 SumMe 与 TVSum 上的总体性能。表中列出的监督方法、无监督方法和 training-free 方法主要用于构成后续正式论文中的横向比较框架；在当前版本中，这些方法的具体数值暂不做主观估计或人工补写，而应以后续查证的原始论文或权威复现结果为准。

\begin{table}[t]
	\centering
	\small
	\caption{第四章方法与代表性视频摘要方法的总体性能对比框架。}\label{tab:c4:overall_performance}
	\begin{tblr}{
		width=\linewidth,
		colspec={Q[l,wd=0.14\linewidth] Q[l,wd=0.30\linewidth] Q[l,wd=0.16\linewidth] Q[c,wd=0.17\linewidth] Q[c,wd=0.17\linewidth]},
		colsep=1pt,
		rowsep=0pt,
		cell{1}{1}={r=2}{c},
		cell{1}{2}={r=2}{c},
		cell{1}{3}={r=2}{c},
		cell{1}{4}={c=2}{c},
	}
		\toprule
		\textbf{类别} & \textbf{方法} & \textbf{监督形式} & \textbf{SumMe Mean $F_1$} $\uparrow$ & \textbf{TVSum Mean $F_1$} $\uparrow$ \\
		\cmidrule{4-5}
		& & & \textbf{结果} & \textbf{结果} \\
		\midrule
		监督方法 & LSTM & Supervised & 待补 & 待补 \\
		监督方法 & dppLSTM & Supervised & 待补 & 待补 \\
		监督方法 & VASNet & Supervised & 待补 & 待补 \\
		监督方法 & PGL-SUM & Supervised & 待补 & 待补 \\
		\midrule
		无监督方法 & KTS & Unsupervised & 待补 & 待补 \\
		无监督方法 & SUM-GAN & Unsupervised & 待补 & 待补 \\
		无监督方法 & DR-DSN & Unsupervised & 待补 & 待补 \\
		无监督方法 & Random / Uniform & Unsupervised & 待补 & 待补 \\
		\midrule
		Training-free 方法 & CLIP-Score & Training-free & 待补 & 待补 \\
		Training-free 方法 & BLIP2-based & Training-free & 待补 & 待补 \\
		Training-free 方法 & LLM-based & Training-free & 待补 & 待补 \\
		Training-free 方法 & Ours (DeepSeek, 8 seg., smoothed / raw-best) & Training-free & 51.89 & 51.41 \\
		Training-free 方法 & Ours (DeepSeek, 16 seg., raw) & Training-free & 待补 & 52.25 \\
		Training-free 方法 & Ours (GPT-4o, 8 seg., smoothed) & Training-free & 待补 & 50.19 \\
		\bottomrule
	\end{tblr}
\end{table}

从当前已经核验的结果看，本文方法在 SumMe 与 TVSum 上均能够稳定得到超过 $50\%$ 的 split 聚合 Mean $F_1$。其中，SumMe 上当前最优的可核验结果来自 DeepSeek 后端、8 段划分并结合平滑后处理的设置，对应 Mean $F_1$ 为 $51.89\%$；TVSum 上当前最优结果来自 DeepSeek 后端、16 段划分且不使用平滑的设置，对应 Mean $F_1$ 为 $52.25\%$。这表明，基于 Planner 与多专家协同推理的摘要框架已经能够在两个标准数据集上形成较稳定的训练外推理能力。

进一步观察当前归档实验可以发现，本文方法在不同设置下的性能变化主要体现在分段粒度、后处理方式以及底层语言模型后端三个方面。对于 SumMe，当前最优结果出现在平滑后处理设置下，说明适度的分数整形有助于提高机器摘要与人工摘要之间的重叠程度；对于 TVSum，更细的分段粒度则表现出更强的摘要筛选能力，说明在视频内容跨度更长、主题转换更频繁的场景中，增加推理单元密度有助于提升关键片段捕获效果。

需要强调的是，表\ref{tab:c4:overall_performance} 中除本文方法外的其他方法目前仅作为正式论文的大表格结构占位，其具体数值仍需在后续定稿阶段结合对应文献逐项补全。因此，现阶段能够得到的严格结论主要限于：一方面，本文方法已经具备可核验的公开基准实验结果；另一方面，其在 SumMe 与 TVSum 上均表现出较稳定的总体性能水平。至于与更多监督、无监督及 training-free 方法之间的细粒度优劣关系，则有待下一小节结合补充完成的文献基线与对照实验进一步分析。

\begin{figure}[t]
	\centering
	\fbox{\parbox[c][4.8cm][c]{0.92\linewidth}{\centering
		图片占位：总体性能柱状图。\\
		建议展示 SumMe 与 TVSum 上各个已完成实验配置的 Mean $F_1$、Precision 和 Recall 对比，可按 DeepSeek-8、GPT-4o-8、DeepSeek-16 等设置分组绘制柱状图，用于体现不同配置下的整体性能水平。
	}}
	\caption{已完成归档实验在公开数据集上的总体性能对比示意图。}\label{fig:c4:overall_performance_bar}
\end{figure}

\begin{figure}[t]
	\centering
	\fbox{\parbox[c][4.8cm][c]{0.92\linewidth}{\centering
		图片占位：5-split 聚合结果分布图。\\
		建议展示 SumMe 与 TVSum 上不同实验配置在 5 个 split 上的 Mean $F_1$ 波动情况，可采用折线图或箱线图，突出不同配置在 split 维度上的稳定性差异。
	}}
	\caption{不同实验配置在 5-split 协议下的性能分布示意图。}\label{fig:c4:overall_performance_split}
\end{figure}

需要指出的是，以上结果均来自当前仓库内已经完成并保存的真实实验记录，而非重新估计或人工补写的数值。与此同时，现有归档结果尚未形成对所有变量的系统扫参，例如 SumMe 上不同后端模型之间的完整对比、TVSum 上更广泛的重叠窗口设置等仍有待补充。因此，本小节的作用主要是给出当前方法的总体性能边界和现阶段最好的可核验结果，而关于不同模型后端、分段粒度与后处理策略之间的更细致差异，下一小节将继续展开分析。

\xsubsection{消融实验}{Ablation Studies}

本小节的目的在于进一步分析第四章方法中若干关键设计选择对摘要性能的影响。与上一小节侧重总体性能汇总不同，这里更关注系统内部配置变化带来的性能差异。需要首先说明的是，依据当前工作区中可直接核验的主执行链路，已经完成并保存结果的可验证消融主要集中在两个层面：其一是分段粒度变化，即按 8 段与 16 段两种方式组织视频推理单元；其二是评测后处理变化，即比较归一化原始分数与归一化平滑分数。对于更细的模块级消融，例如显式关闭记忆、移除 Planner 或裁剪专家集合，当前仓库中虽然保留了部分研究配置痕迹，但尚未在现有主评测链路下形成可直接复核的成套结果，因此本小节对这部分内容仅给出规范化实验骨架，不作虚构性结论。

为便于统一描述，若记某一配置下的 Mean $F_1$ 为 $F_{1}(\mathcal{C})$，则两个配置之间的性能变化可记为：

\begin{equation}
	\Delta F_{1}= F_{1}(\mathcal{C}_{1}) - F_{1}(\mathcal{C}_{2})
\end{equation}
式中，$\Delta F_1$——两种配置之间的 Mean $F_1$ 差值；$F_1(\mathcal{C}_1)$、$F_1(\mathcal{C}_2)$——配置 $\mathcal{C}_1$ 与 $\mathcal{C}_2$ 对应的 Mean $F_1$ 结果；$\mathcal{C}_1$、$\mathcal{C}_2$——待比较的实验设置。

其中 $\mathcal{C}_{1}$ 与 $\mathcal{C}_{2}$ 分别表示两种待比较的实验设置。在当前已落盘结果中，本文主要比较的是真实可运行配置之间的相对差异，而不是对不存在结果的模块做假设性排序。

表\ref{tab:c4:ablation_verified} 汇总了当前 DeepSeek 后端下已经完成的设置消融结果。这里之所以统一固定后端模型，是为了避免将底层语言模型差异与系统内部结构差异混杂在一起。可以看到，8 段与 16 段两种分段粒度，以及 raw 与 smoothed 两种评测变体，已经覆盖了当前主执行链路中最稳定且最容易复现的两类配置变化。

\begin{table}[t]
	\centering
	\small
	\caption{当前主执行链路下已核验的设置消融结果。}\label{tab:c4:ablation_verified}
	\begin{tblr}{
		width=\linewidth,
		colspec={Q[l,wd=0.15\linewidth] Q[l,wd=0.21\linewidth] Q[l,wd=0.26\linewidth] Q[c,wd=0.17\linewidth] Q[c,wd=0.17\linewidth]},
		colsep=1pt,
		rowsep=0pt,
	}
		\toprule
		\textbf{后端} & \textbf{分段设置} & \textbf{评测变体} & \textbf{SumMe Mean $F_1$} $\uparrow$ & \textbf{TVSum Mean $F_1$} $\uparrow$ \\
		\midrule
		DeepSeek & 8 segments & normalized\_raw & 51.41 & 51.41 \\
		DeepSeek & 8 segments & normalized\_smoothed & \textbf{51.89} & 50.73 \\
		DeepSeek & 16 segments & normalized\_raw & 50.66 & \textbf{52.25} \\
		DeepSeek & 16 segments & normalized\_smoothed & 51.12 & 51.81 \\
		\bottomrule
	\end{tblr}
\end{table}

从表\ref{tab:c4:ablation_verified} 可以看出，分段粒度对两个数据集的影响并不完全一致。对于 SumMe，8 段设置整体优于 16 段设置，特别是在 smoothed 变体下，8 段设置取得了 $51.89\%$ 的 Mean $F_1$，而 16 段设置对应结果为 $51.12\%$。这表明在视频长度较短、关键事件相对集中的场景下，过细的分段可能会把原本完整的语义单元切分得过于零散，进而削弱 Planner 对全局主题和片段间关系的稳定判断。

与之相对，TVSum 上的结果呈现出不同趋势。16 段设置在 raw 变体下取得了当前最优的 $52.25\%$ Mean $F_1$，高于 8 段设置的 $51.41\%$。这说明对于主题跨度更大、时长更长的视频集合，更细的时间切分能够为多智能体推理提供更高分辨率的候选片段，使系统更容易定位局部关键事件。不过，当进一步引入平滑后处理时，16 段设置虽然仍保持较高性能，但 Mean $F_1$ 略有回落，说明更高分辨率的片段划分与平滑操作之间并非总是完全协同。

进一步从后处理角度看，平滑操作在两个数据集上的作用也具有明显差异。对于 SumMe，无论是 8 段还是 16 段设置，normalized\_smoothed 的结果都略高于 normalized\_raw，说明在较短视频中对帧级分数进行轻量平滑有助于缓解局部波动，提高机器摘要与人工摘要之间的重叠程度。对于 TVSum，情况则相反：两个分段设置下，smoothed 变体的 Mean $F_1$ 均低于 raw 变体。这说明在内容变化更丰富的视频上，平滑操作虽然可能增强分数曲线的连续性，但也可能削弱部分短时高价值片段的峰值，从而影响最终摘要选择。

需要说明的是，上述结果仍属于设置级消融，其比较对象是当前主评测链路中已经真正运行完成的配置差异，而不是对所有内部模块进行逐一关闭后的严格模块消融。对于模块级消融，当前代码中唯一能够直接从主执行入口确认的显式开关是记忆机制，即通过是否启用 \texttt{enable\_memory} 决定是否向 Planner 与 Expert 提供历史片段上下文。然而，当前工作区中尚未发现与该开关对应的完整数据集级落盘结果，因此无法在本节中给出可靠数值。为保持论文结构完整性，表\ref{tab:c4:ablation_pending} 给出了后续可补充的模块级消融表格骨架。

\begin{table}[t]
	\centering
	\small
	\caption{模块级消融实验结果骨架（待补）。}\label{tab:c4:ablation_pending}
	\begin{tblr}{
		width=\linewidth,
		colspec={Q[l,wd=0.14\linewidth] Q[l,wd=0.34\linewidth] Q[c,wd=0.16\linewidth] Q[c,wd=0.16\linewidth] Q[c,wd=0.12\linewidth]},
		colsep=1pt,
		rowsep=0pt,
	}
		\toprule
		\textbf{变体} & \textbf{设置说明} & \textbf{SumMe Mean $F_1$} $\uparrow$ & \textbf{TVSum Mean $F_1$} $\uparrow$ & \textbf{状态} \\
		\midrule
		Full System & Planner + Experts + Memory + Frame Mapping & 待补 & 待补 & 待补 \\
		w/o Memory & 关闭历史 caption 上下文注入 & 待补 & 待补 & 待补 \\
		w/o Smoothing & 使用 normalized\_raw 评测变体 & 见表\ref{tab:c4:ablation_verified} & 见表\ref{tab:c4:ablation_verified} & 已完成 \\
		16 Segments & 将 \texttt{segment\_value} 调整为 16 & 见表\ref{tab:c4:ablation_verified} & 见表\ref{tab:c4:ablation_verified} & 已完成 \\
		\bottomrule
	\end{tblr}
\end{table}

\begin{figure}[t]
	\centering
	\fbox{\parbox[c][4.8cm][c]{0.92\linewidth}{\centering
		图片占位：已核验消融结果柱状图。\\
		建议展示 DeepSeek 后端下 8 segments / 16 segments 与 raw / smoothed 两两组合在 SumMe 和 TVSum 上的 Mean $F_1$ 对比，以突出分段粒度与后处理方式对性能的不同影响。
	}}
	\caption{当前主执行链路下设置消融结果示意图。}\label{fig:c4:ablation_verified_bar}
\end{figure}

\begin{figure}[t]
	\centering
	\fbox{\parbox[c][4.8cm][c]{0.92\linewidth}{\centering
		图片占位：模块级消融设计示意图。\\
		建议展示 Full System、w/o Memory 以及不同 segment\_value 设置在当前推理流程中的差异位置，标明记忆模块位于 Planner/Expert 打分前的上下文构造阶段，分段粒度位于前端预处理阶段。
	}}
	\caption{模块级与设置级消融位置示意图。}\label{fig:c4:ablation_design}
\end{figure}

综上，当前可以严格核验的消融结果表明，分段粒度与评测后处理对第四章方法的最终摘要性能具有直接影响，而且这种影响在 SumMe 与 TVSum 上表现出明显差异：前者更受益于较粗粒度分段和轻量平滑，后者则更受益于更细粒度的片段建模。与此同时，记忆机制等更贴近方法内部结构的模块级消融仍有待补充完整结果。下一小节将在此基础上，进一步结合具体样例和可视化结果，对方法行为进行更细致分析。

\xsubsection{可视化结果与定性分析}{Visualization and Qualitative Analysis}

为了进一步说明第四章方法在真实视频上的工作方式，本小节结合当前仓库中已经保存的逐视频推理结果与评测结果，从摘要片段选择结果和帧级评分曲线分析两个层面进行定性分析。与上一章基于传统层级语义匹配框架的分析不同，这里关注的是当前多智能体推理链路本身，即 Planner 如何在全局主题约束下调节多专家判断，以及这种段级决策如何映射为最终的帧级摘要分数。需要强调的是，以下分析所涉及的案例均来自现有归档结果中的真实视频，包括 SumMe 中的 ``Cooking''、``Eiffel Tower''，以及 TVSum 中的 ``AwmHb44\_ouw'' 与 ``vdmoEJ5YbrQ''。

\subsubsection{摘要片段选择结果}

从当前归档的 segment 级推理结果看，本文方法在不同视频类型上表现出较一致的选择规律：对于教程类视频，系统会优先保留真正承担步骤推进或关键演示作用的片段；对于景观类与 vlog 类视频，则更强调代表性场景、视角切换以及是否存在明显重复。这一现象与当前框架中的 Planner + Expert 协同机制是一致的，即 Planner 先根据整段视频 caption 建立主题理解，再由 Story、Visual、Emotion 与 Information 四个专家分别提供局部证据，最终形成段级重要性排序。

以 SumMe 中的 ``Cooking'' 为例，当前归档结果将视频前半段大量重复的铁板烹饪操作压到较低分值，而把后半段的火焰展示片段显著抬高。具体而言，segment 1、3、4 的 final score 分别约为 $0.675$、$0.775$ 和 $0.725$，其 Planner reason 均明确指出这些片段与前文在食材和动作上高度重复，新增信息有限；相比之下，segment 6 和 segment 7 的 final score 分别达到 $1.8625$ 与 $1.8375$，对应的 Planner 与 Visual Expert 均将其识别为 ``dramatic flaming presentation''，即整个教程视频的表演性高潮。由此可以看出，当前方法并不是简单偏好动作密集片段，而是能够区分重复的常规操作与具有阶段性收束意义的高光片段，从而使摘要更接近人类在教程视频中的观看偏好。

类似地，在 TVSum 的维修教程视频 ``AwmHb44\_ouw'' 中，系统的高分段主要集中在真正构成维修流程骨架的操作环节。根据归档结果，segment 2、3、4、6 的 final score 分别约为 $1.475$、$1.5625$、$1.35$ 和 $1.4$，对应内容涵盖拆卸螺母、将轮胎从轮辋上撬下、处理轮辋以及后续清理等关键步骤；而仅承担开场介绍或口头说明作用的 segment 0 与 segment 7 得分则分别约为 $1.325$ 和 $1.20$。其中，segment 3 被 Planner 与 Information Expert 同时视为 ``removing the tire from the rim'' 这一核心步骤，因此在视觉冲击并不特别强的情况下仍获得了较高排序。这说明当前多智能体框架在教程类视频中能够优先保留真正具有程序语义和教学价值的内容，而不是仅凭镜头运动或局部纹理变化做出选择。

对于景观类视频 ``Eiffel Tower''，系统的判断标准则明显偏向于是否提供新的代表性视角。当前 8 段归档结果中，segment 3 与 segment 4 的 final score 分别约为 $1.6125$ 与 $1.525$，其高分原因主要来自仰视塔体内部钢结构、穿过拱门观察整体轮廓以及地面游客活动所带来的空间层次变化；而 segment 7 的 final score 下降到 $1.1125$，Planner 明确指出该片段虽然仍然展示铁塔底部结构，但与前面多个低角度镜头相比 ``incremental rather than novel''。这说明对于景观类内容，当前方法不仅关注视觉显著性，还会结合全局上下文抑制重复视角，从而避免摘要被大量相似的地标镜头占满。

再看 TVSum 中的 vlog 视频 ``vdmoEJ5YbrQ''，系统表现出对具有人物互动和环境代表性片段的明显偏好。segment 0 的 final score 为 $1.475$，其高分原因主要在于雪地群体活动、人物出镜与环境全景共同建立了视频的主题基调；segment 1 仍保留较高得分 $1.2625$，因为它补充了人与宠物在雪地中的互动。相对地，视频后段大量车内行驶镜头被显著压低，其中 segment 6 和 segment 7 的 final score 分别仅为 $0.7125$ 与 $0.575$。Planner 在这两段中的判断都强调，其内容虽然维持了冬季出行的连续性，但已经缺乏新的场景信息和互动事件，因此只适合作为过渡而不应大量保留。由此可见，当前方法对 vlog 类视频并不会过度追逐技术上稳定的镜头，而是更重视是否真正承载了旅程中的代表性事件。

\begin{figure}[t]
	\centering
	\fbox{\parbox[c][4.8cm][c]{0.92\linewidth}{\centering
		图片占位：代表性视频的摘要片段选择结果。\\
		建议选取 Cooking、Eiffel Tower、AwmHb44\_ouw 和 vdmoEJ5YbrQ 四个案例，按时间顺序展示原视频片段、系统保留片段以及人工标注高分区域之间的对应关系，并在图中标注高分 segment 编号。
	}}
	\caption{代表性视频的摘要片段选择结果示意图。}\label{fig:c4:qualitative_segment_selection}
\end{figure}

综上，从逐段推理结果可以看到，当前方法在不同视频类型上都能形成相对稳定的保留关键阶段、抑制重复内容的选择逻辑。更重要的是，这种选择不是黑盒地产生，而是能够通过 Planner reason 与各 Expert reason 直接追溯到某一片段为何重要、为何冗余的语义依据，从而使最终摘要具有较强的可解释性。

\subsubsection{帧级评分曲线分析}

在当前实现中，段级最终分数会先被映射到完整帧轴，再根据评测所需的采样位置提取帧级分数。因此，从逐视频评测文件中可以观察到一种非常清晰的曲线结构：原始帧级分数往往呈现明显的分段平台，而平滑后的曲线则在平台边缘形成更加连续的过渡。这一现象本质上来源于第四章方法的实现方式，即片段是当前多智能体推理的基本单位，而帧级结果则是段级语义判断在时间轴上的展开。

以 ``Cooking'' 的评测结果为例，原始 frame scores 在前半段长期维持在中等平台，而在后半段火焰展示对应的时间区域显著抬升。结合推理结果可以发现，这一高平台正对应 segment 6 与 segment 7 两个被 Planner 判定为高潮的片段。由于这两个片段的 final score 明显高于前面的重复烹饪段，映射到帧轴后便表现为连续的高分区间，从而使 knapsack 选择阶段更倾向于保留该部分内容。类似地，在 ``AwmHb44\_ouw'' 中，拆轮胎、撬轮辋和清理等关键步骤所在区域形成了多个相邻高平台，说明教程视频中的程序性强步骤会在帧级曲线上表现为较稳定的连续高分，而非孤立尖峰。

对于 ``Eiffel Tower'' 这类景观视频，帧级曲线则更能体现重复视角抑制的作用。已有评测结果显示，该视频的 frame scores 呈现多个分段常值区间，例如约 $0.7000$、$0.7250$、$1.0000$ 与 $0.8250$ 等平台，其对应的正是不同 segment 在段到帧映射后的展开。由于 segment 3 等片段在 Visual Expert 与 Story Expert 中同时得到较高评价，因此其对应平台被整体抬高；而后段多个重复的塔下仰视镜头虽然仍保留一定分数，但平台高度明显回落。由此可以看出，当前方法对景观视频并不是简单寻找最亮眼的单帧，而是在较长时间尺度上给更具代表性的视角段落以更高权重。

在 vlog 视频 ``vdmoEJ5YbrQ'' 中，这种分段平台特征更加直观。根据当前归档评测结果，其 normalized\_raw frame scores 在视频开头大段时间内保持为 $1.0$ 的高平台，随后依次回落到约 $0.7639$ 和 $0.6111$ 等区间。这一变化与 segment 级推理结论高度一致：前段雪地集体活动和人物互动最能代表视频主题，因此整体得分最高；而后续车内行驶与低质量晃动镜头只承担过渡作用，于是映射后的整段平台被整体下调。也就是说，帧级评分曲线在这里并不是由局部视觉噪声主导，而是基本遵循段级语义重要性的全局排序。

此外，将 normalized\_raw 与 normalized\_smoothed 两种评测变体结合起来看，还可以更清楚地理解上一小节中的消融现象。对于 ``Cooking'' 这类高潮较为集中的视频，平滑操作能够在相邻高分段之间形成更连续的过渡，从而减少局部边界抖动；而对于 TVSum 中某些包含短时关键信息的视频，过度平滑则可能压低原本由段级高分带来的局部峰值，这也与表\ref{tab:c4:ablation_verified} 中 TVSum 更偏好 raw 变体的结果相吻合。因此，帧级曲线不仅可用于解释单个视频的摘要选择，还能够反向说明不同后处理策略为何会对最终 $F_1$ 产生不同影响。

\begin{figure}[t]
	\centering
	\fbox{\parbox[c][4.8cm][c]{0.92\linewidth}{\centering
		图片占位：代表性视频的帧级评分曲线。\\
		建议分别绘制 Cooking、Eiffel Tower、AwmHb44\_ouw 与 vdmoEJ5YbrQ 的 normalized\_raw 和 normalized\_smoothed 曲线，并在下方用色块标出系统最终选择的摘要区间与人工高分区间。
	}}
	\caption{代表性视频的帧级评分曲线与摘要区间示意图。}\label{fig:c4:qualitative_frame_curves}
\end{figure}

总体而言，定性结果与前述定量实验呈现出较强一致性：Planner 负责在全视频层面识别主题与关键阶段，多专家则从叙事推进、视觉高光、情绪强度和信息增量等维度补充局部证据，最终通过段到帧映射形成具有明确结构的分数曲线。正因为这一过程能够同时保留教程视频中的关键步骤、景观视频中的代表性视角以及 vlog 视频中的真实事件节点，本文方法才得以在不同类型数据上保持较稳定的摘要生成能力。

\xsection{本章小结}{summary}

本章围绕开放域长视频在免训练摘要场景下面临的全局理解不足、单一评分器适应性有限以及段级结果难以直接对接标准评测协议等问题，提出了一种面向视频摘要的多智能体大语言模型推理框架。与第三章基于层级语义聚合与相关性匹配的技术路线不同，本章以当前工程主链路为基础，将视频摘要过程组织为视频分段与段级描述生成、Planner 全局规划、多专家协同评分、记忆增强时序建模以及帧级映射与标准评测五个相互衔接的处理环节，从而把原始视频压缩问题转化为可解释的段级协同推理问题。

具体而言，本文首先利用分段机制将长视频重组为若干推理单元，并借助视觉语言模型生成段级多模态描述；在此基础上，由 Planner Agent 从全局范围内推断视频主题、生成全局摘要并动态分配专家权重，再由 Story、Visual、Emotion 与 Information 四类 Expert Agent 对各片段的重要性给出维度化判断；随后，系统通过轻量记忆机制将前序片段语义注入后续评分过程，并将段级最终分数映射到原始采样帧序列，结合归一化、平滑与 knapsack 摘要构造流程生成满足公开数据集协议的最终摘要结果。由此，第四章形成了一个从多模态描述、智能体协同推理到帧级评测输出的完整闭环。

在实验验证方面，当前仓库内已归档结果表明，本章方法能够在 SumMe 与 TVSum 两个标准数据集上形成稳定的公开基准表现。其中，DeepSeek 后端在 8 段划分并采用平滑评测变体时，在 SumMe 上取得了 $51.89\%$ 的已核验 Mean $F_1$；在 16 段划分且采用 raw 变体时，在 TVSum 上取得了 $52.25\%$ 的已核验 Mean $F_1$。进一步的设置级消融说明，分段粒度与后处理策略会对不同数据集产生不同影响：SumMe 更受益于相对较粗的分段和轻量平滑，而 TVSum 更受益于更细粒度的片段建模。结合代表性案例分析还可以看到，当前方法能够在教程类视频中突出关键步骤，在景观类视频中抑制重复视角，并在 vlog 类视频中保留更具事件代表性的片段，从而体现出较好的跨类型适应能力与推理可解释性。

总体来看，本章从当前实际可运行的代码实现出发，完成了多智能体视频摘要框架的建模、实现路径抽象与实验分析，验证了以 Planner 为核心、以多专家协同推理为主体、以帧级映射和标准协议评测为落点的 training-free 视频摘要路线具有可行性。上述工作不仅为本文关于大语言模型驱动视频摘要的研究提供了第二条技术路径，也为下一章对全文工作的整体归纳与进一步讨论奠定了基础。

